\section{Rozpoznawanie kliknięć}
\label{sec:zdarzenia}

\subsection{Pomiar czasu przytrzymania}
Czas przytrzymania mierzą dwa timery z opóźnionym załączeniem (TON),
pobudzane bezpośrednio sygnałem $I_1$:
\begin{itemize}
    \item $T_{001}$ -- ustawiony na \SI{0,5}{\second}; jego wyjście
          przyjmuje wartość 1, gdy przycisk jest wciśnięty co najmniej
          \SI{0,5}{\second},
    \item $T_{002}$ -- ustawiony na \SI{2}{\second}; jego wyjście
          przyjmuje wartość 1, gdy przycisk jest wciśnięty co najmniej
          \SI{2}{\second}.
\end{itemize}
Ponieważ $T_{002}$ może być aktywny tylko wtedy, gdy aktywny jest też
$T_{001}$, zachodzi $T_{002} \Rightarrow T_{001}$.

\subsection{Dlaczego sama detekcja zbocza nie wystarcza}
Rozstrzygnięcie, czy kliknięcie było krótkie, czy średnie, następuje
w chwili puszczenia przycisku. Timery TON zerują się jednak w tym samym
cyklu, w którym przycisk zostaje puszczony. W cyklu zbocza opadającego
wyjścia $T_{001}$ i $T_{002}$ są więc już równe 0 i nie niosą
informacji o czasie trwania naciśnięcia. W pierwszej, nieudanej wersji
programu opartej na impulsie zbocza opadającego i bezpośrednim
odczycie timerów:
\begin{itemize}
    \item kliknięcia krótkie i średnie nie były rozpoznawane,
    \item kliknięcie długie, utrzymywane przez cały czas trzymania
          przycisku, powodowało przełączanie telefonu w każdym cyklu
          (w układzie ze sprzężeniem zwrotnym Set\,$\leftrightarrow$\,Reset),
          czyli szybkie migotanie wyjścia $Q_2$.
\end{itemize}

\subsection{Pamięć poprzedniego cyklu}
Rozwiązaniem jest zapamiętanie wartości sygnałów z \emph{poprzedniego}
cyklu programu. Wykorzystano w tym celu markery $M_7$, $M_8$ i $M_9$,
zapisywane w ostatniej sieci programu. Ponieważ sieć zapisująca marker
jest wykonywana po wszystkich sieciach, które go odczytują,
w sieciach 2--6 marker ma wartość z poprzedniego cyklu:
\begin{equation}
    M_7 := I_1,\qquad M_8 := T_{001},\qquad M_9 := T_{002}.
    \label{eq:poprzedni}
\end{equation}

\subsection{Równania zdarzeń}
Zdarzenia K, Ś i D wyznaczane są jako impulsy trwające jeden cykl:
\begin{align}
    M_4 \;(\mathrm{K}) &= M_7 \,\overline{I_1}\,\overline{M_8}, \label{eq:k}\\
    M_5 \;(\mathrm{Ś}) &= M_7 \,\overline{I_1}\, M_8\,\overline{M_9}, \label{eq:s}\\
    M_6 \;(\mathrm{D}) &= T_{002}\,\overline{M_9}. \label{eq:d}
\end{align}
Interpretacja poszczególnych czynników:
\begin{itemize}
    \item $M_7\,\overline{I_1}$ -- przycisk był wciśnięty w poprzednim
          cyklu, a w bieżącym jest puszczony, czyli wystąpiło zbocze
          opadające,
    \item $\overline{M_8}$ -- w poprzednim cyklu próg \SI{0,5}{\second}
          nie był przekroczony, więc kliknięcie było krótkie,
    \item $M_8\,\overline{M_9}$ -- próg \SI{0,5}{\second} był
          przekroczony, a próg \SI{2}{\second} nie, więc kliknięcie
          było średnie,
    \item $T_{002}\,\overline{M_9}$ -- zbocze narastające $T_{002}$:
          próg \SI{2}{\second} został osiągnięty właśnie w bieżącym
          cyklu, więc wystąpiło kliknięcie długie.
\end{itemize}
Po kliknięciu długim w cyklu puszczenia przycisku $M_8 = M_9 = 1$,
więc ani $M_4$, ani $M_5$ nie przyjmują wartości 1. Spełnia to
wymaganie R8.

\subsection{Tabela prawdy dekodera zdarzeń}
Tabela~\ref{tab:dekoder} przedstawia wartości zdarzeń dla wszystkich
kombinacji $M_7$, $I_1$, $M_8$, $M_9$. Kombinacje z $M_9 = 1$
i $M_8 = 0$ są niemożliwe ($T_{002} \Rightarrow T_{001}$), co
oznaczono jako stan nieokreślony ($-$). Zdarzenie D zależy dodatkowo
od $T_{002}$ i jest pokazane osobno.

\begin{table}[H]
\centering
\caption{Tabela prawdy dekodera kliknięć K i Ś.}
\label{tab:dekoder}
\begin{tabular}{cccc|cc|l}
\toprule
$M_7$ & $I_1$ & $M_8$ & $M_9$ & $M_4$ & $M_5$ & Interpretacja \\
\midrule
0 & $\ast$ & $\ast$ & $\ast$ & 0 & 0 & przycisk nie był wciśnięty \\
1 & 1 & $\ast$ & $\ast$ & 0 & 0 & przycisk nadal wciśnięty \\
1 & 0 & 0 & 0 & \textbf{1} & 0 & puszczony przed \SI{0,5}{\second} -- K \\
1 & 0 & 0 & 1 & $-$ & $-$ & kombinacja niemożliwa \\
1 & 0 & 1 & 0 & 0 & \textbf{1} & puszczony między \SI{0,5}{\second} a \SI{2}{\second} -- Ś \\
1 & 0 & 1 & 1 & 0 & 0 & puszczony po klik. długim -- brak zdarzenia \\
\bottomrule
\end{tabular}
\end{table}

\subsection{Przykładowe przebiegi}
Tabela~\ref{tab:przebieg} przedstawia przebieg kliknięcia średniego
(\SI{1}{\second}) w kolejnych cyklach programu. Impuls $M_5$ trwa
dokładnie jeden cykl.

\begin{table}[H]
\centering
\caption{Przebieg sygnałów dla kliknięcia średniego (\SI{1}{\second}).}
\label{tab:przebieg}
\begin{tabular}{l|ccccc}
\toprule
Moment & $I_1$ & $T_{001}$ & $M_7$ & $M_8$ & $M_5$ \\
\midrule
pierwszy cykl po naciśnięciu    & 1 & 0 & 0 & 0 & 0 \\
$t = \SI{0,3}{\second}$         & 1 & 0 & 1 & 0 & 0 \\
$t = \SI{0,6}{\second}$ (próg)  & 1 & 1 & 1 & 0$\to$1 & 0 \\
$t = \SI{0,9}{\second}$         & 1 & 1 & 1 & 1 & 0 \\
pierwszy cykl po puszczeniu     & 0 & 0 & 1 & 1 & \textbf{1} \\
następny cykl                   & 0 & 0 & 0 & 0 & 0 \\
\bottomrule
\end{tabular}
\end{table}

Dla kliknięcia krótkiego w cyklu puszczenia $M_7=1$ i $M_8=0$, więc
impuls $M_4$ trwa jeden cykl. Dla kliknięcia długiego impuls $M_6$
pojawia się jednokrotnie w cyklu, w którym $T_{002}$ zmienia się
z 0 na 1, a puszczenie przycisku nie generuje zdarzenia.
